\phantomsection
\section*{ВВЕДЕНИЕ}
\addcontentsline{toc}{section}{ВВЕДЕНИЕ}

Нейронный машинный перевод (далее NMT) --- одна из основных задач обработки естественного языка.
С практической точки зрения он заключается в построении модели, которая по тексту на исходном языке формирует текст на целевом языке, сохраняя смысл, грамматическую корректность и контекст высказывания.
Наряду с классическими encoder-decoder системами на основе Transformer все чаще рассматриваются decoder-only языковые модели и архитектуры с линейной обработкой последовательностей.
Результаты последних соревнований WMT показывают, что крупные языковые модели уже стали важной частью систем машинного перевода, однако специализированные NMT-модели по-прежнему остаются конкурентоспособными при контролируемом обучении на параллельных корпусах~\cite{wmt24_general}.

Построение системы NMT включает выбор архитектуры модели, способа представления исходного и целевого текста, схемы организации перевода, набора языковых направлений, состава обучающих данных и метода обучения.
Эти решения совместно определяют качество перевода, вычислительную сложность и возможность переноса модели на новые языки или домены.

С момента появления NMT было предложено несколько основных архитектурных подходов.
К классическим решениям относятся рекуррентные модели LSTM с механизмом внимания~\cite{luong}.
Позднее основой большинства систем стали модели Transformer с механизмом самовнимания~\cite{transformer}.
В последние годы также развиваются архитектуры, основанные на моделях пространства состояний, в частности Mamba~\cite{mamba}.
Эти подходы по-разному работают с контекстом, поэтому их имеет смысл сравнивать не только по качеству перевода, но и по ограничениям при обработке длинных последовательностей.

Не менее важна схема организации самой модели перевода.
В классическом варианте encoder-decoder encoder строит представление исходного предложения, а decoder по этому представлению формирует перевод.
В decoder-only подходе исходный текст и целевой перевод рассматриваются как части одной авторегрессионной последовательности.
Эти схемы могут использоваться с разными архитектурами, поэтому сравнение архитектур без учета способа организации модели дает неполную картину.

Также отличаться может и языковой охват модели.
Модель может обучаться для одного направления перевода, например с одного фиксированного исходного языка на один целевой язык.
Другой вариант --- мультиязычная модель, которая работает сразу с несколькими языками и направлениями перевода.
Мультиязычные модели используют общие представления между языками, но одновременно требуют более сложной балансировки данных, постановки обучения и потенциально другой логики построения перевода.

Современные работы по машинному переводу также показывают, что качество обучающих данных становится не менее важным фактором, чем выбор архитектуры~\cite{newspalm_data}.
Шумные параллельные пары, неверно определенный язык, неполные переводы и технические фрагменты могут ухудшать обучение даже крупных моделей.
Поэтому подготовка корпуса должна рассматриваться как часть общей системы NMT, а не как вспомогательный этап.
Количество техник работы с данными и их комбинаций в переводе велико: фильтрация шумных параллельных предложений, удаление дубликатов, ограничения по длине и соотношению длин предложений, различные способы аугментации данных или генерирование синтетических данных.

Различаются и подходы к обучению моделей перевода.
Базовый вариант --- supervised learning, при котором модель обучается на парах исходных предложений и эталонных переводов.
Дополнительно могут использоваться методы дообучения с учетом автоматических метрик, человеческих оценок или reinforcement learning.
Задача исследования состоит не только в описании отдельных моделей, но и в сопоставлении архитектурных, языковых, корпусных и обучающих решений в общей постановке нейронного машинного перевода.

\subsection*{Постановка задачи}
Целью работы является сравнительное исследование подходов к построению и обучению моделей NMT.
Для достижения поставленной цели необходимо решить следующие задачи:
\begin{enumerate}[label=\arabic*)]
	\item рассмотреть общую постановку задачи NMT;
	\item рассмотреть подходы к подготовке, фильтрации и аугментации обучающих данных;
	\item описать основные подходы к обучению, подбору гиперпараметров и оценке качества моделей перевода;
	\item сравнить архитектуры LSTM, Transformer и Mamba в задачах перевода;
	\item охарактеризовать однонаправленные и мультиязычные модели перевода;
	\item сопоставить encoder-decoder и decoder-only схемы;
\end{enumerate}
